\begin{abstract}
This paper studies regularization for Vision Transformer fine-tuning in a small-data setting. Using ViT-Tiny/16 on CIFAR-100, we build a unified baseline and compare several methods under the same training protocol, including weight decay, dropout-based methods, data augmentation, and label smoothing with early stopping. Our results show clear differences across regularization families. Weight decay provides only small improvements, while dropout-based methods reduce overfitting more effectively. Data augmentation gives the strongest gains among the regularization methods, showing that input-level regularization is especially important for pretrained ViTs in this setting. Label smoothing with early stopping also improves generalization at low implementation cost. In addition, we evaluate layer-wise learning rate decay (LLRD) with staged fine-tuning as an optimization strategy, which further improves performance beyond the baseline.

% \draftnote{Tighten to $\sim$150 words. Follow this 5-sentence skeleton.}
% Sentence 1 (problem): Vision Transformers lack the inductive biases of
% CNNs and tend to overfit on small image datasets such as CIFAR-100.
% Sentence 2 (model + dataset): We fine-tune an ImageNet-pretrained
% ViT-Tiny/16 on CIFAR-100 under a single, fixed training budget so that
% only the regularizer changes between runs.
% Sentence 3 (techniques compared): We study five families of
% regularization---no regularization (baseline), weight decay, dropout,
% data augmentation (Random Crop, Horizontal Flip, Color Jitter,
% RandAugment, Mixup/CutMix), and label smoothing combined with early
% stopping---each swept over a small grid of strengths.
% Sentence 4 (headline finding): \draftnote{One-line headline result, e.g.
% ``Augmentation yields the largest absolute improvement; weight decay is
% the cheapest and most reliable knob; dropout helps only at small rates.''}
% Sentence 5 (conclusion): \draftnote{One-line takeaway about which
% family of regularization matters most for ViT in the small-data
% fine-tuning regime studied here.}
\end{abstract}
